\documentclass[12pt]{article}

\usepackage{times}
\usepackage{ifpdf}
\ifpdf
\usepackage[pdftex]{graphicx}
\else
\usepackage{graphicx}
\fi

\usepackage{minitoc}
\usepackage{subfigure}
\usepackage{color}
\usepackage{url}
\usepackage{latexsym}
\usepackage{geometry}
\usepackage{indentfirst}
\usepackage{array}
\usepackage{tabularx}
\usepackage{xtab}
\usepackage{fancyvrb}
\usepackage[center, bf]{caption}
\usepackage{listings}
\usepackage{url}
\usepackage{placeins}
\usepackage{amsmath}

\begin{document}
\begin{titlepage}
	\centering
	\begin{figure}[ht]
		\centering
		% TODO fix latex errors
		\includegraphics[angle=-90, keepaspectratio=true, width=8cm]{../images/logo}
	\end{figure}
	{
	\large \bfseries Nomadic Comunication\par
	AA 2008-2009
	}\par
	\vspace{1.5cm}
	{
	\Large \bfseries \textcolor{blue}{Report 1: The tile!} \par
	}
	\vspace{1.0cm}
	{
	\large \bfseries {Group N. XXX} \par
	}
	\vspace{0.3cm}
	{
	\large \bfseries {Alberto M. Scattolo, Andrea Crotti, Fabio Cannioto, Christian Plattner}
	}
	\vspace{1.0cm}
	\begin{abstract}
	%Abstract here.
	\end{abstract}
\end{titlepage}

\thispagestyle{empty}
\tableofcontents
\clearpage
\setcounter{page}{1}

%%%%
% Usefull things in latex:
% itemize, enumerate, lstlistings (to report source code, shell output, whatever), emph, tables
%
%%%
% How to include a single figure:
%	\begin{figure}
%		\begin{center}
%			\includegraphics[width=1\textwidth]{path/to/image_file} % use "image_file" for "image_file.pdf"
%			\caption{Figure caption} \label{fig:reference_to_image_file}
%		\end{center}
%	\end{figure}
%
%%%
% TODO
% - Refer to Professor template for some usefull hints: http://disi.unitn.it/locigno/index.php/teaching-duties/nomadic-communications/aa08-09
% - *Always* use \cite and \ref when needed. Adding it when the text is already written is senseless waste of time, the result will be hugly.
% - Use \footnote if necessary, it can make things more readable.
% - Please indent your Latex code! Use tabs, not spaces.
% - English! Write in grammatically correct English. (Aspell can check latex source files)
% - This is a formal document that requires formal writing. This means that "it's", "I've", "ain't" and others must be replaced by their extended form.
% - Use \emph the first time you call something new or when you introduce a new abbreviation i.e.: Internet Protocol (\emph{IP}). Do not abbreviate access point to AP in descriptions and similar
%
%%%%
% IMPORTANT
% Every section is writte in a separate file (see \include commands). This means that anything (subsections and everything else) following \include should be placed in that file. I wrote everything here to make the thing readable.

\section{Introduction} \label{intro}
	% Describe what we are gonna do, the general ideas and what we are gonna talk about.
% Say what is needed to understand the sections that will follow.

The purpose of the project 
is the evaluation of the performances variation of a wireless network 802.11 b/g at the change of some parameters.
\\\newline
In order to make those tests we used some tools which are Iperf, wireshark, and a complex python script for the automated run test.
\\\newline
The test basicaly consists of many automatic executions of iperf with fixed time and different speed, and finaly we consider the amount of transfering data.
\\\newline
The set of the tests was an indoor environment: one notebook which runs iperf as server, another one which runs iperf as client (through the python script), a cisco access point and another notebook wich runs wireshark to capture the traffic.

\vspace{1cm}

\begin{figure}[h!]
	\includegraphics[angle=0, keepaspectratio=true, width=15cm]{../images/network_overview}
	\caption{Network overview}
\end{figure}

	% Describe what we are gonna do, the general ideas and what we are gonna talk about.
	% Say what is needed to understand the sections that will follow.

\section{Theoretical background} \label{theory}
	% Everything that comes from theory and that is needed to understand and describe our work.
	\subsection{802.11} \label{theory:prot_specs}
	802.11 is the standard family of protocols for wireless networks. It is also named Wi-Fi. Protocols in the 802.11 family describe how various wireless devices must interact. There are two kinds of wireless networks: BSS and IBSS.
	\begin{description}
		\item[BSS (infrastructure)] is a centralized network, in this case we have one or more AP (Access Point) and various clients (STAs) can access to the network throw it.
		\item[IBSS (ad-hoc)] is a sample client to client (STA to STA) network.
	\end{description}
	Nowadays the most popular (and commercial) 802.11 implementations are 802.11b and 802.11g.
	
% A simple description of how an IEEE 802.11 protocol works to introduce the fact that there is some overhead added by the protocol.
% Yes, I know it's unbelivable but there are some differences, for real! :)
	\subsection{Differences between 802.11b and 802.11g} \label{theory:prot_differences}
	
	The 802.11b protocol was born in October 1999, and since then it has a discrete success.  However, later on some other wireless devices, such as cordless telephones, use the same frequencies and therefore produce interferences with 802.11b-conforming devices and cause slower performance of these.
	% As a consequence,
	802.11b extends the basic net bit rate of the standard 802.11 up to 11 Mbit/s.\\
	
	In June 2003 arrives 802.11g, which is fully backwards compatible with 802.11b, but brings better bit rates and throughput\\
	
	Both protocols use the same frequency band (2.4 GHz), but different frequency spreads.
	
	\begin{table}[h]
		
		\begin{tabularx}{15cm}{ | X X X X X X | }
			\hline
				802.11 Protocol & Release & Frequency (GHz) & Typical throughput (Mbit/s) & Maximal net bitrate (Mbit/s) & Modulation \\
			\hline
				-- & Jun 1997 & 2.4 & 00.9 & 002 & DSSS \\
				b & Sep 1999 & 2.4 & 04.3 & 011 & DSSS \\
				g & Jun 2003 & 2.4 & 19 & 054 & OFDM \\
			\hline
		\end{tabularx}
		
		\caption{802.11 family standard}
		\label{802.11_family_standard}
	\end{table}
	
	% TODO: ``shy of fivefold then'' ???????
	%As we can see from the table the standard 802.11g's typical throughput is just shy of fivefold then 802.11b standard typical throughput and the main difference is the modulation method. While 802.11b use direct sequence spread spectrum signaling (DSSS), 802.11g use orthogonal frequency division multiplexing (OFDM) methods. Because of OFDM, 802.11g increases the net bit rate up tu 54 Mbit/s.
	%rewrite:
	As we can see from the table, the 802.11g standard's typical throughput is five times that of 802.11b's typical throughput.  The main difference is the modulation method.  While 802.11b uses direct sequence spread spectrum signalling (DSSS), 802.11g uses orthogonal frequency division multiplexing (OFDM) instead.  Because of OFDM, 802.11g increases the net bit rate to up to 54 Mbit/s.
	
	\begin{itemize}
		\item 802.11b (DSSS)
			\begin{description}
				\item[Modulation DBPSK, DQPSK] 1 and 2 Mbit/s
				\item[Modulation CCK] 5.5, 11 Mbit/s
			\end{description}
		\item 802.11g (OFDM)
			\begin{description}
				\item[Modulation BPSK, QPSK, 16QAM, 64QAM] 6, 9, 12, 18, 24, 36, 48, 54 Mbit/s
			\end{description}
	\end{itemize}
%
	% REWRITTEN  When a 802.11g device communicates with a 802.11b device it uses, in order to retrocomputability, DSSS  so it go down with DSSS speeds.
	% TODO check, please:
	When an 802.11g device communicates with an 802.11b device, it switches back to DSSS for backward compatibility. As a consequence, it must communicate with the slower speed of DSSS.

	However, more important than the physical throughput is the effective throughput, which is always lower in relation to the nominal bit rate.
	In 802.11b, we can have a typical throughput of 4,3 Mbit/s over a 11 Mbit/s of channel capacity while in 802.11g we have a typical throuughput of 25 Mbit/s over a 54 Mbit/s of channel capacity. This is due to the overhead of the communication.
	
	The 802.11 standard basically defines two layers, a physical layer (PHY) and a media access control layer (MAC). To understand the reason behind the different throughputs, we have to analyze the MAC Frame Format (Figure \ref{mac_packet}).
	
	\begin{figure}[h!]
		\includegraphics[angle=0, keepaspectratio=true, width=15cm]{../images/mac_header2}
		\caption{MAC Frame Format}
		\label{mac_packet}
	\end{figure}
	
	% TODO: ``need to WEP and MSDU'' ?????
	The total header cost in terms of bytes is fixed to 32 bytes (octets) over a variable size of the frame body from 0 to 2304 bytes plus any overhead from security encapsulation. The frame body size is determined by the maximum MSDU + ICV + IV, where ICV (Integrity Check Value) and IV (Initialization Vector) need to WEP and MSDU is the Mac Service Data Unit. Therefore the real throughput depends by frame body length.\\
	
    % TODO inserti missing figure
    % \begin{figure}[h!]
    %   \includegraphics[angle=0, keepaspectratio=true, width=15cm]{../images/ppdu_frame_format}
    %   \caption{PPDU Frame Format}
    %   \label{ppdu_frame_packet}
    % \end{figure}
	
	To calculate the theoretical throughput with a DATA size of 1470 bit we have to follow the two formulas written below:
	
	\begin{gather*}
		\textrm{Throughput}_{phys} = \\\\
		= \frac{ ( \textrm{DATA}_{bit} + \textrm{HEADER IP}_{bit} + \textrm{HEADER UDP}_{bit} + \textrm{HEADER MAC}_{bit} ) \times 8 }{ \textrm{FRAME}_{seconds} } = \\\\
		= \frac{ 1532 \times 8 }{  \textrm{FRAME}_{seconds} } \\\\
	\end{gather*}
	
	\begin{gather*}
		\textrm{Throughput}_{UDP} = \frac{ \textrm{DATA}_{bit} \times 8 }{ \textrm{FRAME}_{seconds} } = \frac{ 1470 \times 8 }{  \textrm{FRAME}_{seconds} } \\\\
	\end{gather*}
	
	Table \ref{tbl:80211bgexample} shows an example theoretical throughput of the two standards at two different speeds.
	
	\begin{table}[h!]
		\begin{tabularx}{15cm}{ | X | X | X | X | X | }
			\hline
				 & \multicolumn{2}{|c|}{ 802.11b (DSSS)} & \multicolumn{2}{|c|}{ 802.11g-only (OFDM)} \\
				 \hline
				 & 5.5 Mbit/s & 11 Mbit/s & 6 Mbit/s & 54 Mbit/s \\
			\hline
				DIFS & 50 $\mu$s & 50 $\mu$s & 28 $\mu$s & 28 $\mu$s \\
			\hline
				SIFS & 10 $\mu$s & 10 $\mu$s & 10 $\mu$s & 10 $\mu$s \\
			\hline
				CW & 310 $\mu$s & 310 $\mu$s & 139.5 $\mu$s & 139.5 $\mu$s \\
			\hline
				PLPC Preamble + PLPC Header & 96 $\mu$s & 96 $\mu$s & 26 $\mu$s & 26 $\mu$s \\
			\hline
				MAC Header + LLC & 52.36 $\mu$s & 26.18 $\mu$s & 48 $\mu$s & 5.33 $\mu$s \\
			\hline
				IP + UDP Headers & 40.73 $\mu$s & 20.36 $\mu$s & 37.33 $\mu$s & 4.15 $\mu$s \\
			\hline
				Data (1470 byte)& 2138.18 $\mu$s & 1069.09 $\mu$s & 1960 $\mu$s & 217.78 \\
			\hline
				TAIL DATA & -- & -- & 1 $\mu$s & 0.11 $\mu$s \\
			\hline
				ACK PLPC Preamble + Header & 96 $\mu$s & 96 $\mu$s & 26 $\mu$s & 26 $\mu$s \\
			\hline
				ACK & 20.36 $\mu$s & 10.18 $\mu$s & 18.67 $\mu$s & 2.07 $\mu$s \\
			\hline
				TAIL ACK & -- & -- & 1 $\mu$s & 0.25 $\mu$s \\
			\hline
			\hline
				TOTAL & 2813.63 $\mu$s & 1687.81$\mu$s & 2295.5 $\mu$s & 459.19 $\mu$s \\
			\hline
			\hline
				Theoretical physical throughput & 4.35 Mbit/s $\mu$s & 7.26 Mbit/s & 5.33 Mbit/s & 26.69 Mbit/s \\
			\hline
			\hline
				Theoretical physical throughput UDP & 4.17 Mbit/s & 6.96 Mbit/s & 5.12 Mbit/s & 25.61 Mbit/s \\
			\hline
			
		\end{tabularx}
		\caption{Example theoretical throughput for 802.11b and 802.11g}
		\label{tbl:80211bgexample}
	\end{table}
	
	
	\subsection{Fragmentation Threshold and RTS Threshold} \label{theory:frag_rts}
	
	802.11, in order to improve the performance, allows to adjust some parameters like Frag (Fragmentation Threshold) and RTS/CTS (Request To Send / Clear To Send).
	
	\begin{description}
		\item[Fragmentation Threshold] is a technique witch divide a big packet into many small packets. These small packets can be transmitted more quickly into the medium (the air). On the one hand this technique works great avoid the collision but on the other hand it increase the overhead because of the packets header. The possible values range is from 256 to 2346.
	
		\item[RTS Threshold] is the frame size above which an RTS/CTS handshake will be performed before attempting to transmit. The possible values range is from 256 to 2346.
	
		\item[RTS/CTS] is a control signal witch try to synchronize the communication. When a device send an RTS (Request To Send) signal this mean that it is ready to send data, in particular RTS/CTS asks for permission to transmit to reduce collisions, but adds considerable overhead. Disabling RTS/CTS can reduce overhead and latency in WLANs where all stations are close together, but can increase collisions and degrade performance in WLANs where stations are far apart and unable to sense each other to avoid collisions (``Hidden Nodes'', ``Hidden Terminal'').
	\end{description}

% Calculate protocols upperbound limitations and similar things (report the formulas!).
% DO NOT USE PROFESSOR SLIDES AS SOURCE OF INFORMATIONS SINCE MOST OF VALUES ARE COMPLETELY WRONG!!!!!!!!!
% Download the IEEE standard from here: http://disi.unitn.it/locigno/didattica/NC/802.11-2007.pdf and make a large use of ^F.

	









	% Everything that comes from theory and that is needed to understand and describe our work.
	% A simple description of how an IEEE 802.11 protocol works to introduce the fact that there is some overhead added by the protocol.
		%\subsection{Differences between 802.11b and 802.11g} \label{theory:prot_differences}
			% Yes, I know it's unbelivable but there are some differences, for real! :)
			% Calculate protocols upperbound limitations and similar things (report the formulas!).
			% DO NOT USE PROFESSOR SLIDES AS SOURCE OF INFORMATIONS SINCE MOST OF VALUES ARE COMPLETELY WRONG!!!!!!!!!
			% Download the IEEE standard from here: http://disi.unitn.it/locigno/didattica/NC/802.11-2007.pdf and make a large use of ^F.

\section{Testbed Setup} \label{setup}
	% Say what is the testbed used for tests including *everything* necessary, both hardware, software, hardware position, air humidity and everything meaningful for the tests.
	%
	% We have used 2 pc (client C and server S), 1 access pocp reint (AP) and 1 switch. The idea is simple:
	% - AP, C and S are connected to the switch on a wired LAN
	% - we use 2 separate class C LANs (wlan and clan):
	%      one for wired traffic like retriving tcpdump and iperf outputs
	%      the other one for the wireless tests
	% - C connects to S via SSH over clan, starts iperf server and starts tcpdump
	% - C uses iperf to send UDP traffic to S over AP on wlan
	% - C gets out outputs from S using clan
	\subsection{Hardware} \label{setup:hardware}
	% mainly C wifi card configuration, ap settings

\noindent
For the experiment we were provided with some devices by the university, however, we decided to use our own equipment for doing the tests.

Our iperf server and client were run on different notebooks, as shown in table \ref{tbl:laptops}. Furthermore, the network consists of an US Robotics Access Point and a switch from the same company (table \ref{tbl:networkdevices}).


%\subsubsection{Laptops}
%	\begin{tabular*}{1\textwidth}{@{\extracolsep{\fill}} | l l | }
	\begin{table}[h]
		
		\begin{tabularx}{15cm}{ | m{4cm} X | }
			\hline
				ROLE & Client\\
				MODEL & HP Pavilion dv5-1070el\\
				WIFI & Intel PRO/Wireless 5100 AGN\\
				FIRMWERE & iwlwifi-5000 - 6 October 2008\\
			\hline
		\end{tabularx}
		\\\\\\
		\begin{tabularx}{15cm}{ | m{4cm} X | }
			\hline
				ROLE & Server\\
				MODEL & Asus M2410L\\
				WIFI & Intel PRO/Wireless LAN 2100 3B Mini PCI\\
				FIRMWERE & ipw2100\\
			\hline
		\end{tabularx}
		
		\caption{Laptop configuration}
		\label{table1}
		\label{tbl:laptops}
	\end{table}

%\subsubsection{Network devices}
	\begin{table}[h]
		
		\begin{tabularx}{15cm}{ | m{4cm} X | }
			\hline
				ROLE & Access Point\\
				MODEL & US Robotics 805450\\
				FIRMWERE & version 1.53\\
			\hline
		\end{tabularx}
		\\\\\\
		\begin{tabularx}{15cm}{ | m{4cm} X | }
			\hline
				ROLE & Switch\\
				MODEL & US Robotics\\
			\hline
		\end{tabularx}
		
		\caption{Network devices configuration}
		\label{table2}
		\label{tbl:networkdevices}
	\end{table}

	%\subsection{Hardware} \label{setup:hardware}
		% mainly C wifi card configuration, ap settings

	\subsection{Software} \label{setup:software}
	% ssh keys, iwconfig, python software, iperf, tcpdump, gnuplot, whatever used


%For setting up our network we used the iwconfig

\noindent
For measuring the traffic we wrote a python script which makes use of other software.

Python is a simple, interpreted, object oriented scripting language with a large set of libraries.
These are the softwares we used for our project:

\begin{description}
	\item[iperf]
		measures throughput of traffic between two hosts, not limited to wireless setups.
It has a long list of possible options that can be used to test the network with many different low level settings.

The tool can do UDP and TCP traffic tests.\newline
While TCP would measure the throughput which the end-user experiences in most cases, the connectionless UDP option is closer to the low-level details which we were interested in.
To test the network with iperf we need to launch a daemon one host and connect to it.\newline
If we set the interval flag on both the server and the client we get the resulting bandwidht from both, but as we'll see later they could differ.\newline
In fact using the udp mode the client doesn't actually know anything of what is the server actually receiving.
That's why we decided to collect both outputs for every test done.

The only way to understand that the client's results are not precise is to look at the miss rate.
		% more detail?

	\item[tcpdump]
		is a packet sniffer which logs packets sent through one interface. Optionally you can specify a filter on the command line to get only packets you are interested in.
The output file is a binary format, ({\em pcap}), which can be then read by many other applications like Wireshark.
		
	\item[Wireshark]
		is a graphical tool which helps you analyze captured traffic by parsing packets and displaying structured information about them.
		We used this tool in the beginning to see where we can refine our work, but in the end the program was not used anymore.

	\item[gnuplot]
		is an application which can represent data in graphs of many types. It has an interactive text interface, as well as it can be used as a library by a piece of software. We used it with the help of the wrapper library ``python-gnuplot'' for python.

	\item[paramiko]
		is a python library implementing the SSH2 protocol.	
	This is used to enable the python script which automates the tests to launch automatically the iperf server and the tcpdump monitor.

\end{description}
%
The python script was mainly written by Andrea Crotti.  Its main goal is to
automate as many tasks as possible to do our measurement.  The main strength
of the program is its configurability.  From simple configuration ({\em ini})
files it reads the required information about the network topology as well as
the settings to apply when testing the throughput.

A comprehensive description can be found in the file {\tt README.textile}.


\subsubsection{Test automator script}

\noindent
The script \verb|tester.py| is the entry point for the program, and it has support for per-user configuration:
To invoke a test, the program is launched as {\tt ./tester.py \textit{name}},
after placing a configuration file at the location {\tt userconfs/\textit{name}.ini}.\newline

There is also a global configuration file, {\tt default.ini}, which is read at first.
Its values have a lower precedence, thus they are overridden by the per-user configuration file.

For different tests, there are again configuration files inside the {\tt configs} directory, which in turn override some of the configuration before doing the actual experiment.

By default all the configurations get loaded, so every test is a three-merge of default + user\_defined + test\_defined.\newline
This makes a lot easier to define new test sets or test same parameters on different networks.

The contents of these files are described in the next section.


%For the options we considered, see Section \ref{sec:parameters}.


	%\subsection{Software} \label{setup:software}
		% ssh keys, iwconfig, python software, iperf, tcpdump, gnuplot, whatever used

\subsection{Parameters} \label{parameters}
	% Parameters of the tests. Everything that influences the tests and the results we are trying to look at. This includes both static and variables things.
% Wifi channel speed used, fragmentation threshold, rst threshold and similars
\label{sec:parameters}

\noindent
The files are in the INI format, which means they contain simple {\tt\textit{key} = \textit{value}} pairs,
structured into sections by lines with a square-bracketed {\tt[\textit{section-title}]}.

%The script accepts parameters in three different formats:
%\begin{enumerate}
%	\item
%		One literal value, e.\,g.~{\tt speed = 1M},
%
%	\item
%		A {\em list} which will be translated to a list in python,
%		e.\,g.~{\tt speed = 1M, 2M, 5.5M},
%
%	\item
%		A {\em range}, as used in python, like
%		{\tt time = 1..4}, which produces [1, 2, 3, 4].
%\end{enumerate}
%
There are three sections inside the configuration files:

\subsubsection{iperf}

\noindent
This section declares the options passed to the \verb|iperf| tool.
These are the available settings in the section:
% along with the default values we gave them.
% along with the possible options:

\begin{description}
	\item[speed] (iperf argument \verb|-b|, Options: {\tt 1M, 2M, 5.5M, 11M, 36M, 48M, 54M})

		A list of speeds at which the data traffic should be sent over the network.	

		For the test configurations we used only one value each.

	\item[time] (\verb|-t|)

		The time for which the traffic is produced, in seconds.

	\item[interval] (\verb|-i|)

		Sets the time to wait between repeated tests, in seconds.

	\item[format] (\verb|-f|, Default: \verb|K|)

		Tells iperf how it should print data sizes in the result, we only used {\tt K}, which stands for Kilobytes.

	\item[host] (\verb|-c|)

		Defines the server on wich the iperf server will be executed and to which the client will connect.

	\item[udp] (\verb|-u|, Default: True)

		When true, use UDP for measurement, TCP otherwise.
		As mentioned above in \ref{setup:software}, we only use UDP.
\end{description}
%

\subsubsection{ap (Access Point Settings)}

\noindent
The Access Point's settings can not be set automatically, so the script uses this information only for treating the resulting data correctly. Before tests are done, the configuration is printed and the program waits for the user to set everything appropriately and to confirm.

Only the first three of them specify the Access Point's configuration.
All the other options are only stored in the result.
\begin{description}
	\item[mode] (Options: B, G)
	\item[speed] (Options: 1M, 2M, 5.5M, 11M, 36M, 48M, 54M)

		The speed at which the Access Point operates.

	\item[channel] (Options: An integer, usually $[1,13]$)

		The channel used for the network (implies the frequency used).

	\item[ip]	The AP's IP address.  Not used.
	\item[ssid]	The Network's SSID.  Not used.
	\item[firmware]	Information about the driver's firmware.  Not used.
	\item[model]	The Model of the device.  Not used.
	\item[comment]	A specific comment for the user or test.  Not used.
\end{description}
%

\subsubsection{client}

\noindent
Also this section contains mainly informational settings.

\begin{description}
	\item[rts\_threshold] (Threshold for Request-To-Send) %= 256..2437

		Only when the size of a packet is larger than this (in bytes), an RTS will be requested before sending it.
		Otherwise, it will be sent right away.
		This is because for small packet, the RTS becomes a big overhead.

	\item[frag\_threshold] (Fragmentation Threshold) %= 256..1500
		
		Packets will be fragmented and sent piecewise, if their size is bigger than this value.

	\item[brand]	Brand of the network card,
	\item[model]	Model of the network card,
	\item[driver]	Driver of the network card.
\end{description}
%

\subsubsection{monitor}

\noindent
This section defines where and how \verb|tcpdump| will be executed.
\begin{description}
	\item[interface]
		Determines on which network interface \verb|tcpdump| will listen.

	\item[host]
		Specifies the monitor's IP or DNS address.
		On this host, an SSH server must be running, and root's public key must be authorized on the local machine.
\end{description}


	% Parameters of the tests. Everything that influences the tests and the results we are trying to look at. This includes both static and variables things.
	% Wifi channel speed used, fragmentation threshold, rst threshold and similars

\section{Results} \label{results}
	The tests we executed are in order
%(andrea/tests)s

\begin{figure}
    
\end{figure}
	% Our results and something that underlines that we understand the mean of the results. Results are not self evident!
	% Any formula used to calculate the efficency or whatever must be written and described.
	%\subsection{IEEE 802.11b} \label{resutls:11b}
	%\subsection{IEEE 802.11g} \label{resutls:11g}

\section{Conclusion} \label{concluion}
	% Conclusion must say what we understand from our results, what is the learned lesson and so on.

\addcontentsline{toc}{section}{Bibliography}
\bibliographystyle{plain}
\bibliography{biblio} % put bibliografy in a file called biblio.bib
% Include all books, manuals, softwares home pages etc used for the tests.
% Including links to wikipedia is like saying that we are a bunch of fools... better not to :)
\nocite{*}
\appendix % if necessary
% We should do something for Andrea's software...
\end{document}
